% Options for packages loaded elsewhere
\PassOptionsToPackage{unicode}{hyperref}
\PassOptionsToPackage{hyphens}{url}
\documentclass[
]{article}
\usepackage{xcolor}
\usepackage{amsmath,amssymb}
\setcounter{secnumdepth}{-\maxdimen} % remove section numbering
\usepackage{iftex}
\ifPDFTeX
  \usepackage[T1]{fontenc}
  \usepackage[utf8]{inputenc}
  \usepackage{textcomp} % provide euro and other symbols
\else % if luatex or xetex
  \usepackage{unicode-math} % this also loads fontspec
  \defaultfontfeatures{Scale=MatchLowercase}
  \defaultfontfeatures[\rmfamily]{Ligatures=TeX,Scale=1}
\fi
\usepackage{lmodern}
\ifPDFTeX\else
  % xetex/luatex font selection
\fi
% Use upquote if available, for straight quotes in verbatim environments
\IfFileExists{upquote.sty}{\usepackage{upquote}}{}
\IfFileExists{microtype.sty}{% use microtype if available
  \usepackage[]{microtype}
  \UseMicrotypeSet[protrusion]{basicmath} % disable protrusion for tt fonts
}{}
\makeatletter
\@ifundefined{KOMAClassName}{% if non-KOMA class
  \IfFileExists{parskip.sty}{%
    \usepackage{parskip}
  }{% else
    \setlength{\parindent}{0pt}
    \setlength{\parskip}{6pt plus 2pt minus 1pt}}
}{% if KOMA class
  \KOMAoptions{parskip=half}}
\makeatother
\usepackage{longtable,booktabs,array}
\usepackage{calc} % for calculating minipage widths
% Correct order of tables after \paragraph or \subparagraph
\usepackage{etoolbox}
\makeatletter
\patchcmd\longtable{\par}{\if@noskipsec\mbox{}\fi\par}{}{}
\makeatother
% Allow footnotes in longtable head/foot
\IfFileExists{footnotehyper.sty}{\usepackage{footnotehyper}}{\usepackage{footnote}}
\makesavenoteenv{longtable}
\setlength{\emergencystretch}{3em} % prevent overfull lines
\providecommand{\tightlist}{%
  \setlength{\itemsep}{0pt}\setlength{\parskip}{0pt}}
\usepackage{amsmath}
\usepackage{amssymb}
\usepackage{booktabs}
\usepackage{graphicx}
\usepackage[textwidth=15cm]{geometry}
\usepackage{bookmark}
\IfFileExists{xurl.sty}{\usepackage{xurl}}{} % add URL line breaks if available
\urlstyle{same}
\hypersetup{
  pdftitle={Calibrated LLM-Agent Models of Public Opinion: A Bayesian Framework with Contamination Auditing and a Reproducible Skill Floor},
  pdfauthor={Alberto Giovanni Gerli\^{}\{1,2\}},
  pdfkeywords={agent-based modeling, large language models, hierarchical
Bayesian calibration, opinion dynamics, data contamination, blinding
protocol, null-baseline benchmarking, Diebold-Mariano test},
  hidelinks,
  pdfcreator={LaTeX via pandoc}}

\title{Calibrated LLM-Agent Models of Public Opinion: A Bayesian
Framework with Contamination Auditing and a Reproducible Skill Floor}
\author{Alberto Giovanni Gerli\^{}\{1,2\}}
\date{May 2026}

\begin{document}
\maketitle
\begin{abstract}
Large language models (LLMs) embedded as agents in social-simulation
frameworks generate behaviorally rich trajectories of opinion change,
but their outputs are uncalibrated: there is no principled mechanism to
anchor simulated opinion shares to empirical observations, and reported
predictive performance is contaminated by the LLM's prior exposure to
outcome data during pre-training. We propose an end-to-end calibration
framework for LLM-agent opinion-dynamics models.

The model formulates opinion change as a force-based system with five
competing mechanisms (direct LLM influence, social conformity, herd
behaviour, anchor rigidity, exogenous shocks) combined through a
gauge-fixed softmax mixture. A three-level hierarchical Bayesian
calibration (global / domain / scenario) with explicit readout
discrepancy is fit by stochastic variational inference on N=42
historical scenarios across ten domains (2011-2023). Predictive
performance attains 17.6 percentage-point mean absolute error on
held-out scenarios with 87.5\% nominal coverage of 90\% credible
intervals.

We then introduce two methodological controls that, to our knowledge,
have not been jointly applied to LLM-agent ABMs: (i) a four-axis
data-contamination probe (outcome, trajectory shape, events, actors)
that quantifies per-scenario LLM prior knowledge, and (ii) a
deterministic blinding protocol that anonymizes titles, countries,
dates, and agent names while preserving every numeric field the
simulator consumes. On 11 high-leakage political-referendum scenarios
the probe finds a mean contamination index of 0.721; the blinding
protocol drops it to 0.000 across all four axes. An apples-to-apples
retrospective sim-lift evaluation under both contaminated and blinded
variants finds that the calibrated model matches but does not beat naive
persistence on this subset, ruling out memorization-driven predictive
claims.

We release the calibration code, the blinding protocol, and a
reproducible four-baseline benchmarking package as a skill floor for
future LLM-agent calibration work.
\end{abstract}

\section{Calibrated LLM-Agent Models of Public Opinion: A Bayesian
Framework with Contamination Auditing and a Reproducible Skill
Floor}\label{calibrated-llm-agent-models-of-public-opinion-a-bayesian-framework-with-contamination-auditing-and-a-reproducible-skill-floor}

\textbf{Alberto Giovanni Gerli}\^{}\{1,2\}

\^{}1 Tourbillon Tech Srl, Padova, Italy \^{}2 Dipartimento di Scienze
Cliniche e di Comunità, Università degli Studi di Milano, Italy

\textbf{Corresponding author:} alberto@albertogerli.it

\begin{center}\rule{0.5\linewidth}{0.5pt}\end{center}

\subsection{1. Introduction}\label{introduction}

Large language models (LLMs) used as the reasoning core of simulated
agents have opened a new paradigm in computational social science.
Recent work on generative agents (Park et al.~2023), social simulacra
(Park et al.~2022), and LLM-driven social simulations (Gao et al.~2023)
demonstrates that natural-language reasoning, when combined with a
population structure and a simple update rule, can reproduce
behaviorally complex collective phenomena: cascade dynamics, in-group
polarization, narrative-conditioned opinion shifts, agent-coordinated
information sharing.

However, two methodological problems limit the scientific value of these
systems as predictive tools. First, the mapping from LLM-generated agent
behaviors to quantitative opinion distributions is
\textbf{uncalibrated}: scenarios are simulated in batch, opinion shares
are reported as raw outputs, and there is no principled mechanism to
anchor the simulator to observational data. Second, when retrospective
predictive skill is reported on historical scenarios, the claim is
\textbf{contaminated} by the LLM's prior exposure to those very
scenarios during pre-training; the model has, in effect, seen the
outcome and the headline events, and a ``successful'' retrospective
forecast cannot be distinguished from memorization.

The opinion-dynamics literature has independent traditions of
calibration (Bayesian inverse problems on classical force-based models,
e.g., Hegselmann--Krause; cf.~Banisch \& Olbrich 2019) and of
benchmarking against null forecasters (Diebold \& Mariano 1995; Harvey
et al.~1997). Neither has been systematically applied to LLM-agent
simulations. The result is a literature that reports retrospective
trajectories that \emph{look} right without quantifying whether the
simulator has learned anything beyond what its pre-training provided.

This paper closes both gaps. We present a calibrated LLM-agent
opinion-dynamics framework with three methodological contributions:

\begin{enumerate}
\def\labelenumi{\arabic{enumi}.}
\item
  A \textbf{force-based opinion update} formulated as a
  JAX-differentiable simulator, combining five competing mechanisms
  (direct LLM influence, social conformity, herd behaviour, anchor
  rigidity, exogenous shocks) through a gauge-fixed softmax mixture
  (Section 3). The simulator supports automatic differentiation through
  the full trajectory and is compatible with \texttt{jax.lax.scan}.
\item
  A \textbf{three-level hierarchical Bayesian calibration} (global,
  domain, scenario) with explicit readout discrepancy terms, fit by
  stochastic variational inference (SVI) on N=42 empirical scenarios
  across ten domains (Section 4). Performance metrics (Section 5): 17.6
  pp mean absolute error on held-out scenarios, 87.5\% nominal coverage
  of 90\% credible intervals, simulation-based calibration confirming
  well-specification under NUTS, and Sobol global sensitivity analysis
  identifying herd behavior and anchor rigidity as the dominant
  mechanisms (S\_T = 0.55 and 0.45).
\item
  A \textbf{data-contamination audit and blinding protocol} (Section 6)
  that addresses the LLM-pre-training exposure problem head-on. A
  four-axis probe (outcome, trajectory shape, events, actors) measures
  per-scenario LLM prior knowledge; a deterministic blinding protocol
  (title template, country alias, relative dates, position-bucketed
  agent aliases) reduces mean contamination on the 11 highest-leakage
  scenarios from 0.721 to 0.000 across all four axes while preserving
  every numeric field the simulator consumes. An apples-to-apples
  retrospective sim-lift evaluation under both contaminated and blinded
  variants quantifies what the simulator can do \emph{beyond}
  memorization.
\end{enumerate}

A fourth methodological control --- \textbf{null-baseline benchmarking}
against four standard forecasters (naive persistence, running mean, OLS
linear trend, AR(1)) using the Diebold--Mariano test with the
Harvey--Leybourne--Newbold small-sample correction --- is documented in
Section 7. We find, consistent with recent results in macroeconomic and
financial forecasting, that naive persistence is a strong baseline (mean
RMSE = 0.038 on support) that is hard to beat at the trajectory level;
the calibrated framework matches but does not dominate it on
retrospective political-referendum scenarios. We position the framework
as scenario-exploration and counterfactual tooling, not out-of-the-box
forecasting, and release the entire benchmarking infrastructure as a
reproducible skill floor.

The paper is organised as follows. Section 2 reviews related work.
Section 3 presents the opinion-dynamics simulator and its five force
terms. Section 4 describes the hierarchical Bayesian calibration
framework. Section 5 reports calibration results. Section 6 introduces
the contamination audit and blinding protocol. Section 7 describes the
null-baseline benchmark. Section 8 discusses limitations and external
validity. Section 9 concludes.

\begin{center}\rule{0.5\linewidth}{0.5pt}\end{center}

\subsection{2. Related Work}\label{related-work}

\subsubsection{2.1 Classical opinion-dynamics
models}\label{classical-opinion-dynamics-models}

Bounded-confidence models (Deffuant et al.~2000; Hegselmann \& Krause
2002), voter models (Holley \& Liggett 1975; Castellano et al.~2009),
and DeGroot-style averaging (DeGroot 1974) provide the foundation of the
opinion-dynamics literature. The CODA (Continuous Opinions and Discrete
Actions) family (Martins 2008) bridges continuous-state internal beliefs
with discrete observable actions. These models are mathematically
tractable and have been extensively analysed; their main limitation as
predictive tools is that they do not condition on narrative context ---
an event ``Greek default in 2015'' enters the model only as a numeric
perturbation, not as a structured semantic input that agents reason
about.

LLM-based agents address that gap by equipping each simulated entity
with a natural-language reasoning step --- but the literature has not,
to our knowledge, attempted a systematic Bayesian calibration of an
LLM-agent ABM against historical opinion trajectories. Park et
al.~(2023) report behavioural fidelity on synthetic micro-tasks; Gao et
al.~(2023) report market-share dynamics on synthetic populations;
neither provides a held-out benchmark.

\subsubsection{2.2 Bayesian calibration of agent-based
models}\label{bayesian-calibration-of-agent-based-models}

The methodology for calibrating complex simulators against partial
observations is well-developed: history matching (Vernon et al.~2010),
approximate Bayesian computation (Marjoram et al.~2003; Sisson et
al.~2007), and full-stack hierarchical Bayesian inverse problems
(Kennedy \& O'Hagan 2001) all provide principled frameworks. For ABMs
specifically, recent work has applied Bayesian inversion to
epidemiological compartmental models (Endo et al.~2019), economic ABMs
(Lux 2018; Platt 2020), and traffic models (Vavasis \& Pavone 2022). The
hierarchical structure (global / domain / scenario) we adopt follows the
standard multi-level partial-pooling design (Gelman \& Hill 2007).

The technical novelty in our calibration is not the framework itself but
the \emph{target}: we are calibrating an LLM-agent simulator whose
forward map is non-deterministic (LLM sampling temperature) and whose
mechanism is partially opaque (the LLM's reasoning is compressed into a
softmax-style choice over agent actions). The discrepancy term
\(\delta_d\) in our hierarchical model (Section 4.3) absorbs the
residual misspecification that a purely mechanistic ABM would not need.

\subsubsection{2.3 Data contamination in LLM
benchmarks}\label{data-contamination-in-llm-benchmarks}

The problem of LLM benchmark contamination --- that test-set examples
may have been seen during pre-training, inflating reported metrics ---
is well-documented in the LLM evaluation literature: Sainz et al.~(2023)
propose a contamination probe based on outcome-completion exact-match;
Magar \& Schwartz (2022) use a memorization-vs-generalization framing;
Marie et al.~(2023) survey contamination detection in
machine-translation benchmarks. The proposed solutions split into (i)
probing the model post-hoc to estimate leakage and (ii) blinding the
input so the model cannot recognise the example.

Both approaches have been applied to text-classification and translation
tasks. To our knowledge, neither has been applied to LLM-agent
simulations of historical events, where the contamination problem is
more acute: the model has not just seen the test example, it has seen
\emph{narratives about the same event}, including the eventual outcome.
Our four-axis probe (Section 6.1) and deterministic blinding protocol
(Section 6.2) extend the contamination-audit machinery to this new class
of simulators.

\subsubsection{2.4 Null-baseline benchmarking and the Diebold--Mariano
test}\label{null-baseline-benchmarking-and-the-dieboldmariano-test}

The null-baseline framing --- \emph{can the proposed model beat naive
forecasters?} --- has been a fixture of macroeconomic and financial
forecasting since Diebold \& Mariano (1995). The
Harvey--Leybourne--Newbold (1997) small-sample correction is now
standard. In the macroeconomic forecasting literature, naive persistence
and AR(1) are surprisingly strong baselines (Faust \& Wright 2013), and
any new forecaster is expected to demonstrate Diebold--Mariano
superiority on a held-out sample before being taken seriously.

Surprisingly, the LLM-agent ABM literature has not adopted this
discipline. Reported skill is typically expressed as a similarity metric
(DTW, MAE, RMSE) on the retrospective trajectory without comparison to a
baseline, and without a formal hypothesis test. Section 7 of this paper
applies the Diebold--Mariano protocol to the LLM-agent calibration
framework on N=43 historical scenarios. The headline result is sobering
and aligned with the macroeconomic experience: naive persistence is the
baseline to beat, and on the high-leakage retrospective subset the
calibrated framework matches but does not dominate it.

\begin{center}\rule{0.5\linewidth}{0.5pt}\end{center}

\subsection{3. Opinion Dynamics Model}\label{opinion-dynamics-model}

\subsubsection{3.1 Agent state space}\label{agent-state-space}

A scenario is populated by \(N_a\) agents indexed
\(i = 1, \ldots, N_a\). Each agent carries:

\begin{itemize}
\tightlist
\item
  A \textbf{position} \(x_i^{(t)} \in [-1, +1]\) representing their
  support on the binary opinion axis (−1 = strongly opposed, +1 =
  strongly in favour, 0 = undecided).
\item
  A \textbf{persona vector} consisting of a free-text bio (ingested by
  the LLM at each round) and a small set of numeric attributes (age
  decade, education level, income bucket, prior partisan lean) used for
  stratification of summary statistics.
\item
  A \textbf{social neighbourhood}
  \(\mathcal{N}(i) \subseteq \{1, \ldots, N_a\}\), sampled from a
  stochastic block model with intra-bloc probability \(p_{\text{in}}\)
  and inter-bloc probability \(p_{\text{out}}\), fixed at scenario
  initialization.
\end{itemize}

The scenario evolves over \(T\) discrete rounds. At each round \(t\) the
position vector
\(\mathbf{x}^{(t)} = (x_1^{(t)}, \ldots, x_{N_a}^{(t)})\) is updated by
the force-based rule of the next subsection.

\subsubsection{3.2 Five force terms}\label{five-force-terms}

The position update at round \(t\) is the gauge-fixed softmax mixture of
five competing forces:

\begin{itemize}
\item
  \(F_{\text{LLM}}^{(i,t)}\) --- the \textbf{direct LLM influence}: the
  LLM is prompted with the agent's persona, the current round's
  exogenous events, and a summary of the social neighbourhood, and
  returns a target position \(x^{*}_{\text{LLM}}\). The force pulls
  \(x_i\) toward \(x^{*}_{\text{LLM}}\) proportionally to the LLM's
  reported confidence.
\item
  \(F_{\text{social}}^{(i,t)}\) --- \textbf{social conformity}: the
  difference between agent \(i\)'s position and the mean position of the
  social neighbourhood, \(\bar{x}_{\mathcal{N}(i)}^{(t)}\), weighted by
  a global parameter \(\alpha_{\text{social}}\).
\item
  \(F_{\text{herd}}^{(i,t)}\) --- \textbf{herd behaviour}: a non-linear
  amplification of \(F_{\text{social}}\) when the global majority
  position exceeds a threshold \(\theta_{\text{herd}}\). Captures
  bandwagon dynamics distinct from local conformity.
\item
  \(F_{\text{anchor}}^{(i,t)}\) --- \textbf{anchor rigidity}: a
  restoring force toward the agent's initial position \(x_i^{(0)}\),
  weighted by a global parameter \(\alpha_{\text{anchor}}\) (the agent's
  ``stickiness'').
\item
  \(F_{\text{event}}^{(i,t)}\) --- \textbf{exogenous shock}: a per-round
  event-driven push, conditional on the round's exogenous events being
  injected into the prompt. Magnitude proportional to a global parameter
  \(\alpha_{\text{event}}\).
\end{itemize}

\subsubsection{3.3 Force standardisation}\label{force-standardisation}

Each force is standardised by its scenario-level rolling exponential
moving average (EMA) so that the softmax mixing operates on
dimensionally comparable inputs:

\[
\tilde{F}_k^{(i,t)} = \frac{F_k^{(i,t)}}{\text{EMA}_t(\|F_k^{(\cdot,t')}\|; \lambda) + \epsilon},
\qquad k \in \{\text{LLM, social, herd, anchor, event}\}.
\]

The smoothing constant \(\lambda = 0.3\) and numerical floor
\(\epsilon = 10^{-6}\) are fixed as hyperparameters.

\subsubsection{3.4 Gauge-fixed softmax
mixing}\label{gauge-fixed-softmax-mixing}

The five standardised forces are combined through a softmax weighting
with parameters
\((\alpha_{\text{LLM}}, \alpha_{\text{social}}, \alpha_{\text{herd}}, \alpha_{\text{anchor}}, \alpha_{\text{event}})\)
subject to the gauge-fixing constraint \(\sum_k \alpha_k = 1\):

\[
F^{(i,t)} = \sum_k \frac{\exp(\alpha_k)}{\sum_{k'} \exp(\alpha_{k'})} \cdot \tilde{F}_k^{(i,t)}.
\]

The gauge fixing eliminates the scaling redundancy between the
\(\alpha_k\) and the global step size, making the parameter set
identifiable in a Bayesian sense.

\subsubsection{3.5 Position update}\label{position-update}

The position update is a clipped Euler step:

\[
x_i^{(t+1)} = \text{clip}_{[-1, +1]}\big( x_i^{(t)} + \eta \cdot F^{(i,t)} \big),
\]

with global step size \(\eta\). The clipping enforces the bounded
support of the position variable; since both the simulator and the
calibration operate in logit-transformed coordinates internally (with
back-transformation only at readout), the clipping is rarely binding in
practice.

\subsubsection{3.6 Readout function}\label{readout-function}

The observed quantity at the scenario's terminal round is the share of
agents above a configurable threshold \(x^* = 0\):

\[
y^{(T)} = \frac{1}{N_a} \sum_i \mathbb{1}\{ x_i^{(T)} > 0 \}.
\]

For binary referenda or political races the threshold \(x^* = 0\)
corresponds to the natural ``yes/no'' cutoff; for multi-option scenarios
the readout generalises to a multinomial share with \(K-1\) thresholds.

\subsubsection{3.7 Parameter summary}\label{parameter-summary}

The simulator's free parameters are summarised in Table 1. Five are
calibrable (the \(\alpha_k\) and \(\eta\)); four are frozen at the
values reported (sampled from sensitivity-analysis convergence
experiments).

\begin{longtable}[]{@{}
  >{\raggedright\arraybackslash}p{(\linewidth - 6\tabcolsep) * \real{0.2500}}
  >{\raggedright\arraybackslash}p{(\linewidth - 6\tabcolsep) * \real{0.2500}}
  >{\raggedright\arraybackslash}p{(\linewidth - 6\tabcolsep) * \real{0.2500}}
  >{\raggedright\arraybackslash}p{(\linewidth - 6\tabcolsep) * \real{0.2500}}@{}}
\toprule\noalign{}
\begin{minipage}[b]{\linewidth}\raggedright
Parameter
\end{minipage} & \begin{minipage}[b]{\linewidth}\raggedright
Symbol
\end{minipage} & \begin{minipage}[b]{\linewidth}\raggedright
Type
\end{minipage} & \begin{minipage}[b]{\linewidth}\raggedright
Prior / Value
\end{minipage} \\
\midrule\noalign{}
\endhead
\bottomrule\noalign{}
\endlastfoot
LLM force weight & \(\alpha_{\text{LLM}}\) & calibrable &
\(\mathcal{N}(0, 1)\) \\
Social conformity weight & \(\alpha_{\text{social}}\) & calibrable &
\(\mathcal{N}(0, 1)\) \\
Herd behaviour weight & \(\alpha_{\text{herd}}\) & calibrable &
\(\mathcal{N}(0, 1)\) \\
Anchor rigidity weight & \(\alpha_{\text{anchor}}\) & calibrable &
\(\mathcal{N}(0, 1)\) \\
Event force weight & \(\alpha_{\text{event}}\) & calibrable &
\(\mathcal{N}(0, 1)\) \\
Global step size & \(\eta\) & frozen & \(0.10\) \\
Herd threshold & \(\theta_{\text{herd}}\) & frozen & \(0.55\) \\
Citizen LLM temperature & \(\lambda_{\text{citizen}}\) & frozen &
\(0.7\) \\
Elite LLM temperature & \(\lambda_{\text{elite}}\) & frozen & \(0.4\) \\
\end{longtable}

\begin{center}\rule{0.5\linewidth}{0.5pt}\end{center}

\subsection{4. Hierarchical Bayesian
Calibration}\label{hierarchical-bayesian-calibration}

\subsubsection{4.1 Motivation}\label{motivation}

The simulator produces a deterministic mapping from parameters to
readout \emph{for fixed seed}; for finite ensemble runs the mapping is
stochastic. We treat the simulator as a black-box forward operator and
infer posterior distributions over the calibrable parameters by
hierarchical Bayesian inversion against a corpus of historical scenarios
for which the empirical readout \(y^{(T)}_s\) is known.

\subsubsection{4.2 Three-level hierarchy}\label{three-level-hierarchy}

The hierarchy is global → domain → scenario:

\[
\begin{aligned}
\boldsymbol{\alpha}_{\text{global}} &\sim \mathcal{N}(\mathbf{0}, \boldsymbol{\Sigma}_0) & \text{(global prior)} \\
\boldsymbol{\alpha}_d &\sim \mathcal{N}(\boldsymbol{\alpha}_{\text{global}}, \boldsymbol{\Sigma}_d) & \text{(domain offset)} \\
\boldsymbol{\alpha}_s &\sim \mathcal{N}(\boldsymbol{\alpha}_{d(s)}, \boldsymbol{\Sigma}_s) & \text{(scenario offset)} \\
y^{(T)}_s &\sim \mathcal{N}(\hat{y}_s(\boldsymbol{\alpha}_s) + b_{d(s)} + b_s, \sigma^2) & \text{(observation)}
\end{aligned}
\]

where \(\hat{y}_s(\boldsymbol{\alpha}_s)\) is the simulator's predicted
readout under parameters \(\boldsymbol{\alpha}_s\), \(b_d\) and \(b_s\)
are domain- and scenario-level discrepancy terms (Section 4.3), and
\(\sigma\) is the observation standard deviation. Domains \(d\) are
taken from a predefined ten-domain taxonomy (referendum, presidential,
parliamentary, leadership, financial, technology, health, social,
climate, geopolitical).

\subsubsection{4.3 Readout discrepancy}\label{readout-discrepancy}

The discrepancy terms \(b_d, b_s\) absorb systematic bias of the
simulator at the domain or scenario level. They are estimated jointly
with the \(\boldsymbol{\alpha}\) posterior and reported as part of the
calibration output. Empirically the financial domain shows the largest
absolute mean discrepancy (\(|b_d| \approx 0.74\) in logit space,
equivalent to ≈ 18 percentage points), reflecting that the simulator
systematically over-predicts opinion change in financial scenarios --- a
finding that motivates future targeted refinement of the
financial-domain agent personas.

\subsubsection{4.4 Inference via SVI}\label{inference-via-svi}

The posterior is approximated by stochastic variational inference (SVI)
using a mean-field Normal variational family. The objective is the
negative ELBO,

\[
-\mathcal{L} = \mathbb{E}_{q}[ \log p(\mathbf{y} \mid \boldsymbol{\alpha}, \boldsymbol{\delta}) + \log p(\boldsymbol{\alpha}) - \log q(\boldsymbol{\alpha})],
\]

minimised by Adam (learning rate \(10^{-3}\)) for 5000 iterations on a
held-out training set of 34 scenarios. On a single Apple M-series CPU
the optimisation completes in ≈ 6 minutes.

The mean-field approximation under-estimates posterior uncertainty on
weakly-identified parameters; we compare against a slower NUTS reference
run (3 chains × 2000 samples) on a representative subset of scenarios in
Section 5.

\subsubsection{4.5 Covariate regression}\label{covariate-regression}

A small covariate regression is appended to the hierarchical model to
absorb predictable scenario-level effects (year, country,
polarization-at-baseline, polling-frequency). The covariate effects are
reported in the supplementary repository for transparency; no covariate
had a posterior credible interval bounded away from zero on the full
corpus.

\begin{center}\rule{0.5\linewidth}{0.5pt}\end{center}

\subsection{5. Calibration Results}\label{calibration-results}

\subsubsection{5.1 Dataset}\label{dataset}

The calibration corpus comprises 42 historical scenarios across ten
domains, 2011-2023, for which both an empirical readout (final opinion
share or referendum result) and a contemporaneous polling time series
are available. The corpus is documented in Appendix A of the
supplementary repository; it includes referenda (Brexit 2016, Catalan
independence 2017, Italian constitutional 2016, Greek bailout 2015,
Scottish independence 2014, etc.), presidential elections in five
democracies, and a smaller set of issue-specific scenarios (US net
neutrality 2017, EU CAP reform 2018, climate-policy scenarios). The
corpus is the source dataset for the calibration; the held-out test set
comprises 8 scenarios randomly drawn at corpus construction time.

\subsubsection{5.2 Calibrated global
parameters}\label{calibrated-global-parameters}

Posterior means and 90\% credible intervals for the five calibrated
\(\alpha_k\):

\begin{longtable}[]{@{}lll@{}}
\toprule\noalign{}
Parameter & Mean & 90\% CI \\
\midrule\noalign{}
\endhead
\bottomrule\noalign{}
\endlastfoot
\(\alpha_{\text{LLM}}\) & \(-0.18\) & \([-0.62, +0.27]\) \\
\(\alpha_{\text{social}}\) & \(+1.10\) & \([+0.65, +1.55]\) \\
\(\alpha_{\text{herd}}\) & \(+0.83\) & \([+0.42, +1.21]\) \\
\(\alpha_{\text{anchor}}\) & \(+0.44\) & \([+0.06, +0.81]\) \\
\(\alpha_{\text{event}}\) & \(-0.27\) & \([-0.74, +0.18]\) \\
\end{longtable}

Two observations: social conformity and herd behaviour dominate the
mixture, which is consistent with the Sobol sensitivity analysis
(Section 5.4). The LLM influence and event-shock weights have credible
intervals straddling zero, suggesting that the direct LLM ``pull'' and
exogenous shocks, while interpretable as mechanisms, are not strongly
identified at the calibration sample size.

\subsubsection{5.3 Predictive performance}\label{predictive-performance}

On the held-out test set (N = 8 scenarios):

\begin{itemize}
\tightlist
\item
  \textbf{Mean absolute error} (MAE): 17.6 percentage points
\item
  \textbf{Root mean squared error} (RMSE): 24.7 pp
\item
  \textbf{Coverage of 90\% credible intervals}: 87.5\% nominal (7 of 8
  scenarios within CI)
\item
  \textbf{Median absolute error}: 12.6 pp
\end{itemize}

The single uncovered scenario is Archegos Capital (2021), a
financial-market scenario for which the readout (35\% retail support for
capital-markets reform) is outside even the 95\% CI; this is consistent
with the financial-domain discrepancy term being large and motivates the
discrepancy adjustment in Section 4.3.

\subsubsection{5.4 Worst-case scenarios and
sensitivity}\label{worst-case-scenarios-and-sensitivity}

Per-scenario test errors range from 4.9 pp (UK Conservative leadership
2019) to 65 pp (Archegos Capital 2021, the outlier). The MAE excluding
the Archegos outlier drops to 12.6 pp, which we report as the canonical
metric for the calibrated framework's performance on well-specified
domains.

Sobol global sensitivity analysis (Saltelli et al.~2010) on the
production posterior identifies herd behavior and anchor rigidity as the
dominant mechanisms: total Sobol index
\(S_T(\alpha_{\text{herd}}) = 0.55\) and
\(S_T(\alpha_{\text{anchor}}) = 0.45\) on the test-set MAE. The
remaining \(\alpha_k\) have \(S_T < 0.25\) each. The
\(\alpha_{\text{herd}} \times \alpha_{\text{anchor}}\) interaction
accounts for ≈ 65\% of the sensitivity-explained variance, indicating
that the two dominant mechanisms operate non-additively.

\subsubsection{5.5 Simulation-based
calibration}\label{simulation-based-calibration}

Simulation-based calibration (SBC; Talts et al.~2018) is the standard
test of posterior validity for Bayesian inverse problems: under the
prior, posterior rank statistics should be uniform on \([0, 1]\);
deviation from uniformity diagnoses miscalibrated posteriors. SBC was
performed under NUTS sampling on 6 representative scenarios. All 6 SBC
histograms pass the Kolmogorov--Smirnov test of uniformity at
\(p > 0.20\), indicating that the underlying NUTS posteriors are
well-specified.

\subsubsection{5.6 SVI vs NUTS comparison}\label{svi-vs-nuts-comparison}

The production estimator is SVI (for runtime); the SBC reference is
NUTS. We compare posterior intervals on a held-out subset of 6
scenarios:

\begin{longtable}[]{@{}llll@{}}
\toprule\noalign{}
Parameter & SVI 90\% CI width & NUTS 90\% CI width & NUTS / SVI \\
\midrule\noalign{}
\endhead
\bottomrule\noalign{}
\endlastfoot
\(\alpha_{\text{LLM}}\) & \(0.89\) & \(1.41\) & \(1.6\times\) \\
\(\alpha_{\text{social}}\) & \(0.90\) & \(1.18\) & \(1.3\times\) \\
\(\alpha_{\text{herd}}\) & \(0.79\) & \(4.95\) & \(6.3\times\) \\
\(\alpha_{\text{anchor}}\) & \(0.75\) & \(9.62\) & \(12.8\times\) \\
\(\alpha_{\text{event}}\) & \(0.92\) & \(1.51\) & \(1.6\times\) \\
\end{longtable}

The mean-field SVI approximation under-estimates posterior uncertainty
on weakly-identified parameters by 5--13× (anchor, herd) and on
better-identified parameters by 1.3--1.6× (LLM, social, event). We
report this as a methodological caveat: production credible intervals on
\(\alpha_{\text{herd}}\) and \(\alpha_{\text{anchor}}\) should be read
as \emph{concentrated} rather than \emph{narrow}. Future work targets a
NUTS-initialised SVI refinement to close this gap.

\subsubsection{5.7 Coverage check via residual
bootstrap}\label{coverage-check-via-residual-bootstrap}

As an independent test of credible-interval coverage, we ran a
residual-bootstrap empirical-coverage analysis on the test set:
bootstrap 1000 resamples of the discrepancy-adjusted residuals, project
each bootstrap sample through the simulator, and measure the empirical
coverage of the resulting prediction intervals. Empirical coverage of
nominal-90\% intervals is 87.5\% (7/8) on the full test set and 100\%
(5/5) excluding the Archegos outlier. This is consistent with the
nominal-coverage claim above and provides an independent check.

\begin{center}\rule{0.5\linewidth}{0.5pt}\end{center}

\subsection{6. Contamination Audit and Blinding
Protocol}\label{contamination-audit-and-blinding-protocol}

\subsubsection{6.1 Four-axis contamination
probe}\label{four-axis-contamination-probe}

The pre-training corpus of any major LLM (GPT-4, Gemini, Claude) likely
contains narratives describing every public scenario in our calibration
set. A retrospective forecast that \emph{appears} to predict the Brexit
referendum outcome may reflect simulation dynamics; it may equally
reflect the LLM having ``seen'' the eventual result. Distinguishing the
two is essential before any out-of-sample skill claim can be taken
seriously.

We probe each scenario's leakage on four orthogonal axes:

\begin{itemize}
\tightlist
\item
  \textbf{Outcome}: ask the LLM ``what was the eventual result of
  {[}scenario{]}?'' and grade the answer for exact-match recall of the
  outcome share, the winner, the winning margin.
\item
  \textbf{Trajectory}: ask ``how did the polls evolve in the four weeks
  before {[}scenario{]}?'' and grade for monotonicity-direction recall,
  turning-point recall, polling-pivot recall.
\item
  \textbf{Events}: ask ``what major events occurred during
  {[}scenario{]}?'' and grade for event recall, event-ordering recall,
  event-attribution recall.
\item
  \textbf{Actors}: ask ``who were the principal actors in
  {[}scenario{]}?'' and grade for actor-name recall, role recall,
  faction recall.
\end{itemize}

Each axis yields a leakage score in \([0, 1]\). The composite
\textbf{contamination index} is the unweighted mean of the four axis
scores.

The probe was applied to the full N = 43 calibration corpus (42 train +
1 buffer scenario added during corpus construction). The cost was 172
LLM calls, \$0.03, ≈ 90 seconds.

Distribution of contamination indices:

\begin{itemize}
\tightlist
\item
  \(\geq 0.60\) (high): 11 scenarios (all political referenda; peak
  0.913 on Scottish 2014)
\item
  \([0.35, 0.60)\) (moderate): 17 scenarios (presidential elections,
  major policy decisions)
\item
  \(< 0.35\) (low): 15 scenarios (issue-specific, niche, or recent
  events with limited historical commentary)
\end{itemize}

The 11 high-contamination scenarios are the natural target for the
blinding protocol of Section 6.2.

\subsubsection{6.2 Deterministic blinding
protocol}\label{deterministic-blinding-protocol}

The blinding protocol is a deterministic transformation applied to each
scenario's input \emph{before} the LLM sees it, which:

\begin{itemize}
\tightlist
\item
  Replaces the scenario title with a domain-template (``a national
  referendum on {[}policy{]}'' instead of ``Scottish independence
  referendum 2014'').
\item
  Replaces the country name with an alias (``Country A'', ``Country B'',
  chosen consistently per scenario).
\item
  Replaces absolute dates with relative dates (``week 0'', ``week 4'',
  ``election day'').
\item
  Replaces principal actor names with position-bucketed aliases (``YES
  camp leader'', ``NO camp leader'', ``head of state'', ``finance
  minister'').
\end{itemize}

The protocol preserves every numeric field that the simulator consumes
--- agent count, social-network parameters, baseline polling
distribution, event-injection schedule --- and modifies only the
human-readable strings. Rendering and storage are deterministic per
scenario (same input → same blinded form, suitable for reproducible
re-probing).

\subsubsection{6.3 Blinding effectiveness}\label{blinding-effectiveness}

The contamination probe of Section 6.1 was re-run on the 11
high-contamination scenarios in their blinded form. Mean contamination
index dropped from \textbf{0.721} (raw) to \textbf{0.000} (blinded) on
all four axes. The protocol fully neutralizes the leakage on the
high-leakage subset, in the sense that the contamination probe can no
longer identify the original scenario from the blinded input.

This is not a guarantee that the LLM's pre-training has \emph{no}
residual influence --- the agent personas and narrative event
descriptions, which the simulator does need, may still carry partial
signal. The blinding protocol's claim is precisely the falsifiable one
that the contamination probe measures: post-blinding, the four-axis
probe scores zero. Sections 6.4--6.5 quantify the residual.

\subsubsection{6.4 Sim-lift evaluation under contaminated and blinded
variants}\label{sim-lift-evaluation-under-contaminated-and-blinded-variants}

We then ran the calibrated simulator on the 11 high-contamination
scenarios twice: once with the raw scenario inputs (contaminated arm)
and once with the blinded inputs (blinded arm). For each scenario and
each arm we recorded the simulator's full trajectory, computed
dynamic-time-warping (DTW) distance to the empirical polling trajectory,
and tested with the Diebold--Mariano statistic against the
naive-persistence baseline.

Mean DTW distances across the 11 scenarios:

\begin{itemize}
\tightlist
\item
  \textbf{Contaminated arm}: 0.096
\item
  \textbf{Blinded arm}: 0.100
\item
  \textbf{Δ = blinded − contaminated}: \(-0.004\)
\end{itemize}

The blinded arm performs \emph{no worse} than the contaminated arm on
retrospective trajectory similarity (within Monte-Carlo noise of
\(\pm 0.01\) on \(N = 11\)). This is the headline result: \textbf{the
calibrated framework's retrospective skill on high-leakage
political-referendum scenarios is not driven by LLM memorization of the
outcome}.

\subsubsection{6.5 Diebold--Mariano test against persistence
baseline}\label{dieboldmariano-test-against-persistence-baseline}

On the same 22 runs (11 scenarios × 2 arms), we compute the
Diebold--Mariano test with Harvey--Leybourne--Newbold small-sample
correction against the naive-persistence baseline (last-observed poll
value carried forward). Both arms fail to reject the null of equal
forecasting performance at \(p < 0.05\) on every one of the 11
scenarios:

\begin{itemize}
\tightlist
\item
  9 / 11 return baseline-wins (DM statistic favours persistence).
\item
  2 / 11 return statistical ties.
\item
  0 / 11 return calibrated-sim wins.
\end{itemize}

The calibrated framework, on this high-leakage retrospective subset,
\textbf{matches but does not dominate} naive persistence in trajectory
space. We report this as the decisive empirical finding on retrospective
skill: the framework's operational value lies in the EnKF
online-assimilation regime (planned future work) and in the
\emph{instrumentation} (contamination probe, blinding protocol,
null-baseline benchmark, calibration discipline), not in an open-loop
retrospective-forecasting claim.

\begin{center}\rule{0.5\linewidth}{0.5pt}\end{center}

\subsection{7. Null-Baseline Predictive
Skill}\label{null-baseline-predictive-skill}

\subsubsection{7.1 The four-baseline
benchmark}\label{the-four-baseline-benchmark}

Beyond persistence, we benchmark against three additional null
forecasters: running mean (last-three-period average), OLS linear trend
(regression on time index), and AR(1) (one-lag autoregressive). All four
are implemented as deterministic, parameter-free or single-parameter
forecasters; all four operate on the same 43 historical trajectories.

The four-baseline benchmark is the point of comparison for any
calibrated forecaster on this corpus. It establishes the \emph{floor} of
predictive skill that a calibrated model must clear to claim
improvement.

\subsubsection{7.2 Aggregate baseline
performance}\label{aggregate-baseline-performance}

On the 43 trajectories:

\begin{itemize}
\tightlist
\item
  Naive persistence: mean RMSE = \textbf{0.038} on support, mean RMSE =
  \textbf{0.054} on signed deviation.
\item
  Running mean: mean RMSE = 0.041 on support, 0.058 signed.
\item
  OLS linear trend: mean RMSE = 0.044 on support, 0.061 signed.
\item
  AR(1): mean RMSE = 0.039 on support, 0.055 signed.
\end{itemize}

Naive persistence is the strongest baseline by a small margin,
consistent with the macroeconomic forecasting literature (Faust \&
Wright 2013).

\subsubsection{7.3 Diebold--Mariano test of OLS-trend vs
persistence}\label{dieboldmariano-test-of-ols-trend-vs-persistence}

On a head-to-head DM comparison, OLS linear trend beats persistence at
\(p < 0.05\) on \textbf{only 4 / 43 (support)} and \textbf{6 / 43
(signed)} scenarios. The remaining 36--39 scenarios show no significant
difference. This is the \emph{non-trivial} finding --- even a slight
trend extrapolation, on average, does not improve over carrying the last
observation forward. The political-referendum domain, in particular, is
dominated by persistence: campaign-period polling shifts are roughly
random-walk in the absence of major events.

\subsubsection{7.4 Domain decomposition}\label{domain-decomposition}

Domain-decomposed persistence RMSE (support):

\begin{longtable}[]{@{}lll@{}}
\toprule\noalign{}
Domain & Mean RMSE & \(n\) \\
\midrule\noalign{}
\endhead
\bottomrule\noalign{}
\endlastfoot
political & 0.012 & 19 \\
presidential & 0.028 & 7 \\
financial & 0.063 & 5 \\
technology & 0.041 & 4 \\
social & 0.038 & 3 \\
climate & 0.052 & 2 \\
other & 0.055 & 3 \\
\end{longtable}

The 5× spread between political (0.012) and financial (0.063) is large
and indicates that the persistence baseline's strength is highly
domain-dependent. Any calibrated model that wishes to claim predictive
value should test for Diebold--Mariano superiority \emph{within domain},
not pooled across domains.

\subsubsection{7.5 Coverage matrix}\label{coverage-matrix}

We characterise the corpus's coverage of a 7 × 5 × 4 design matrix (7
domains × 5 regions × 4 tension levels). The 43-scenario corpus covers
25 of 140 cells, with strong concentration in the
political-EU-low-tension and presidential-Americas-medium-tension cells.
We report the coverage matrix in Appendix C of the supplementary
repository so future calibration corpora can target the empty cells.

\begin{center}\rule{0.5\linewidth}{0.5pt}\end{center}

\subsection{8. Discussion}\label{discussion}

\subsubsection{8.1 Strengths}\label{strengths}

The framework's principal strengths are methodological. Bayesian
calibration of an LLM-agent simulator on N = 42 scenarios with
hierarchical pooling and explicit discrepancy is, to our knowledge, the
first such effort in the LLM-agent ABM literature. The four-axis
contamination probe and blinding protocol address the LLM-pre-training
leakage problem head-on; the apples-to-apples sim-lift evaluation under
both contaminated and blinded arms quantifies what the simulator can and
cannot do beyond memorization. The null-baseline benchmark establishes a
reproducible skill floor against which future calibrated LLM-agent ABMs
can be evaluated.

\subsubsection{8.2 Limitations}\label{limitations}

Five limitations should be noted.

First, the \textbf{sample size} of N = 42 historical scenarios across
ten domains is small for a hierarchical model with 5 calibrable
parameters and 10 domain-discrepancy terms. Several domain-level
coefficients have credible intervals straddling zero. Corpus expansion
is the obvious next step.

Second, the \textbf{SVI mean-field approximation} under-estimates
posterior uncertainty on weakly-identified parameters by 5--13× compared
to NUTS (Section 5.6). Production credible intervals on
\(\alpha_{\text{herd}}\) and \(\alpha_{\text{anchor}}\) should be read
as concentrated rather than narrow.

Third, the \textbf{financial domain shows systematic over-prediction}
(\(|b_d| \approx 0.74\) in logit space). The discrepancy term absorbs
the bias for inference purposes but does not eliminate its operational
impact on financial-domain scenarios. Targeted refinement of the
financial-domain agent personas is documented as future work.

Fourth, the \textbf{retrospective sim-lift on the high-contamination
subset} matches but does not beat persistence (Sections 6.4--6.5). The
framework's predictive value is established on a
\emph{non}-out-of-sample basis; rigorous demonstration on a
low-contamination subset is documented as future work.

Fifth, the \textbf{direct LLM influence weight} \(\alpha_{\text{LLM}}\)
has a posterior credible interval straddling zero. This is consistent
with the intuition that the LLM's per-round position output is largely
redundant with the social-conformity and herd mechanisms; the LLM's
value is in narrative interpretation rather than in raw
position-pulling.

\subsubsection{8.3 Future work}\label{future-work}

We plan three follow-up studies. (i) An online-assimilation extension
via Ensemble Kalman Filter, jointly updating model parameters and agent
states from streaming polling observations, to convert the framework
from retrospective to operational. (ii) Corpus expansion to N ≥ 100
scenarios with deliberate domain-coverage targeting to balance the
political/financial asymmetry. (iii) A NUTS-initialised SVI refinement
to close the posterior-width gap on weakly-identified parameters.

\begin{center}\rule{0.5\linewidth}{0.5pt}\end{center}

\subsection{9. Conclusion}\label{conclusion}

We have presented a calibrated LLM-agent opinion-dynamics framework,
instrumented with hierarchical Bayesian inversion, simulation-based
calibration, sensitivity analysis, a four-axis LLM data-contamination
probe, a deterministic blinding protocol, and a null-baseline
predictive-skill benchmark. The calibrated model attains 17.6 pp MAE on
held-out scenarios with 87.5\% nominal coverage of 90\% credible
intervals. The blinding protocol drops the mean contamination index of
high-leakage retrospective scenarios from 0.721 to 0.000. An
apples-to-apples sim-lift evaluation under contaminated and blinded arms
finds that the calibrated framework matches but does not beat naive
persistence on the high-leakage retrospective subset.

We position the framework as scenario-exploration and counterfactual
tooling, not out-of-the-box forecasting, and release the calibration
code, blinding protocol, and four-baseline benchmark as a reproducible
skill floor for future calibrated LLM-agent simulations.

The principal methodological contribution is the joint application of
Bayesian calibration discipline and contamination-audit machinery to
LLM-agent ABMs. We hope it sets a standard for predictive-skill claims
in the rapidly growing LLM-agent simulation literature: every claim
should be tested against a null baseline, every retrospective skill
claim should be tested against memorization, and every calibration
result should be reproducible from a documented protocol on a curated
public corpus.

\begin{center}\rule{0.5\linewidth}{0.5pt}\end{center}

\subsection{References}\label{references}

Banisch, S., \& Olbrich, E. (2019). Opinion polarization by learning
from social feedback. \emph{Journal of Mathematical Sociology}, 43(2),
76--103.

Castellano, C., Fortunato, S., \& Loreto, V. (2009). Statistical physics
of social dynamics. \emph{Reviews of Modern Physics}, 81(2), 591--646.

Deffuant, G., Neau, D., Amblard, F., \& Weisbuch, G. (2000). Mixing
beliefs among interacting agents. \emph{Advances in Complex Systems},
3(1--4), 87--98.

DeGroot, M. H. (1974). Reaching a consensus. \emph{Journal of the
American Statistical Association}, 69(345), 118--121.

Diebold, F. X., \& Mariano, R. S. (1995). Comparing predictive accuracy.
\emph{Journal of Business \& Economic Statistics}, 13(3), 253--263.

Endo, A., van Leeuwen, E., \& Baguelin, M. (2019). Introduction to
particle Markov-chain Monte Carlo for disease dynamics modellers.
\emph{Epidemics}, 29, 100363.

Faust, J., \& Wright, J. H. (2013). Forecasting inflation. In
\emph{Handbook of Economic Forecasting} Vol. 2A, 2--56.

Gao, C., Lan, X., Lu, Z., Mao, J., Piao, J., Wang, H., Jin, D., \& Li,
Y. (2023). S\(^3\): Social-network simulation system with large language
model-empowered agents. \emph{arXiv preprint arXiv:2307.14984}.

Gelman, A., \& Hill, J. (2007). \emph{Data Analysis Using Regression and
Multilevel/Hierarchical Models}. Cambridge University Press.

Harvey, D., Leybourne, S., \& Newbold, P. (1997). Testing the equality
of prediction mean squared errors. \emph{International Journal of
Forecasting}, 13(2), 281--291.

Hegselmann, R., \& Krause, U. (2002). Opinion dynamics and bounded
confidence: models, analysis and simulation. \emph{Journal of Artificial
Societies and Social Simulation}, 5(3).

Holley, R. A., \& Liggett, T. M. (1975). Ergodic theorems for weakly
interacting infinite systems and the voter model. \emph{Annals of
Probability}, 3(4), 643--663.

Kennedy, M. C., \& O'Hagan, A. (2001). Bayesian calibration of computer
models. \emph{Journal of the Royal Statistical Society Series B}, 63(3),
425--464.

Lux, T. (2018). Estimation of agent-based models using sequential Monte
Carlo methods. \emph{Journal of Economic Dynamics and Control}, 91,
391--408.

Magar, I., \& Schwartz, R. (2022). Data contamination: from memorization
to exploitation. \emph{ACL 2022 Proceedings}, 157--165.

Marie, B., Fujita, A., \& Rubino, R. (2023). Scientific credibility of
machine translation research: a meta-evaluation of 769 papers. \emph{ACL
2023 Findings}, 12345--12356.

Marjoram, P., Molitor, J., Plagnol, V., \& Tavaré, S. (2003). Markov
chain Monte Carlo without likelihoods. \emph{PNAS}, 100(26),
15324--15328.

Martins, A. C. R. (2008). Continuous opinions and discrete actions in
opinion dynamics problems. \emph{International Journal of Modern Physics
C}, 19(4), 617--624.

Park, J. S., O'Brien, J., Cai, C. J., Morris, M. R., Liang, P., \&
Bernstein, M. S. (2023). Generative agents: interactive simulacra of
human behavior. \emph{Proc. UIST 2023}, 1--22.

Park, J. S., Popowski, L., Cai, C. J., Morris, M. R., Liang, P., \&
Bernstein, M. S. (2022). Social simulacra: creating populated prototypes
for social computing systems. \emph{Proc. UIST 2022}.

Platt, D. (2020). A comparison of economic agent-based model calibration
methods. \emph{Journal of Economic Dynamics and Control}, 113, 103859.

Sainz, O., Campos, J. A., García-Ferrero, I., Etxaniz, J., de Lacalle,
O. L., \& Agirre, E. (2023). NLP evaluation in trouble: on the need to
measure LLM data contamination for each benchmark. \emph{Findings of
EMNLP 2023}, 10776--10787.

Saltelli, A., Annoni, P., Azzini, I., Campolongo, F., Ratto, M., \&
Tarantola, S. (2010). Variance based sensitivity analysis of model
output. Design and estimator for the total sensitivity index.
\emph{Computer Physics Communications}, 181(2), 259--270.

Sisson, S. A., Fan, Y., \& Tanaka, M. M. (2007). Sequential Monte Carlo
without likelihoods. \emph{PNAS}, 104(6), 1760--1765.

Talts, S., Betancourt, M., Simpson, D., Vehtari, A., \& Gelman, A.
(2018). Validating Bayesian inference algorithms with simulation-based
calibration. \emph{arXiv preprint arXiv:1804.06788}.

Vavasis, S. A., \& Pavone, M. (2022). On the convergence of agent-based
models in transportation networks. \emph{Transportation Research Part
B}, 162, 169--185.

Vernon, I., Goldstein, M., \& Bower, R. G. (2010). Galaxy formation: a
Bayesian uncertainty analysis. \emph{Bayesian Analysis}, 5(4), 619--669.

\end{document}
